% !PRINTED none
% Verlagsanzeige von Friedrich Vieweg & Sohn in Braunschweig; ohne Seitenzahl. Auf den
% Verlagskopf und einen kurzen Mittelstrich folgen fünf vollständige, durch kurze
% Mittelstriche getrennte Anzeigen:
% 1. „Praktische Anleitung zur Anstellung astronomischer Beobachtungen mit besonderer
% Rücksicht auf die Astrophysik. Nebst einer modernen Instrumentenkunde“ von Nicolaus von
% Konkoly, Dr. phil., Mitglied der Akademie der Wissenschaften in Budapest, der Royal
% Astronomical Society in London etc.; mit 345 Holzstichen; gr. 8.; geh.; Preis 24 Mark.
% 2. „Vorträge und Reden“ von Hermann von Helmholtz; vierte Auflage; mit dem Bildniss des
% Verfassers und zahlreichen eingedruckten Holzstichen; zwei Bände; gr. 8.; pro Band geh.
% 8 Mark, geb. 9 Mark 50 Pfennig.
% 3. „Der Mond und die Beschaffenheit und Gestaltung seiner Oberfläche“ von Edmund Neison,
% Mitglied der Königlichen astronomischen Gesellschaft zu London etc.; autorisirte deutsche
% Original-Ausgabe; zweite Auflage, vermehrt mit einem Anhange, enthaltend die Untersuchungen
% des Verfassers über die Neubildung Hyginus N auf dem Monde; nebst einem Atlas in 4° von
% 26 Karten und 5 Tafeln in Farbendruck; gr. 8.; geh.; mit Atlas zusammen 18 Mark.
% 4. „Photogrammetrie und internationale Wolkenmessung“ von Dr. Carl Koppe, Professor an
% der Herzogl. technischen Hochschule zu Braunschweig; mit Abbildungen und fünf Tafeln;
% gr. 8.; geh.; Preis 7 Mark.
% 5. „Fünfstellige logarithmische und trigonometrische Tafeln“, herausgegeben von
% Dr. O. Schlömilch, K. S. Geheimerath a. D., Mitglied der Königl. Sächsischen Gesellschaft
% der Wissenschaften zu Leipzig, der Königl. Schwedischen Akademie zu Stockholm, der
% Kaiserl. Leopoldinischen Akademie etc.; galvanoplastische Stereotypie; wohlfeile
% Schulausgabe; zwölfte Auflage; 8.; geh.; Preis 1 Mark. Die Anzeige endet auf dieser Seite.
